% Solutions for chapter 10-hochschild
\providecommand{\cA}{\mathcal{A}}
\providecommand{\cB}{\mathcal{B}}

\section{Solutions for \cref{ch:hochschild}: Hochschild homology}\label{solutions:hochschild}

%----------------------------------------------------------------------
\subsection*{\Cref{hochschild:ex:diagonal}}
Idea: a projective $\cA^e$-module is its own projective resolution, of length zero; the
comparison theorem of \cref{hochschild:sec:Tor} then computes $\HH_\bullet$.

Let $e_1,\dots,e_N$ be the standard idempotents of $\cA=\C^N$, so $e_ie_j=\delta_{ij}e_i$ and
$1=\sum_je_j$. In $\cA^e=\cA\otimes\cA^{\mathrm{op}}$ (with product $(a\otimes a')(c\otimes c')=ac\otimes c'a'$) the
element $\epsilon=\sum_je_j\otimes e_j$ satisfies
$\epsilon^2=\sum_{i,j}e_ie_j\otimes e_je_i=\sum_je_j\otimes e_j=\epsilon$. Let $m:\cA^e\to\cA$,
$m(a\otimes a')=aa'$, be the multiplication map (the augmentation $P_0\to\cA$ of
\eqref{hochschild:eq:bar-res}), and define
\begin{equation*}
  s:\cA\longrightarrow\cA^e,\qquad s(a)=(a\otimes1)\epsilon=\sum_jae_j\otimes e_j .
\end{equation*}
Then $m(s(a))=\sum_jae_j=a$. Moreover $s$ is $\cA^e$-linear: for $b,b'\in\cA$, since $b'e_j=b'_je_j$
($b'_j$ the $j$-th component),
\begin{equation*}
  (b\otimes b')\cdot s(a)=\sum_jbae_j\otimes e_jb'=\sum_jb'_j\,bae_j\otimes e_j=\sum_jbab'e_j\otimes e_j=s(bab')=s((b\otimes b')\cdot a).
\end{equation*}
So $m\circ s=\Id$ and $\cA^e=s(\cA)\oplus\ker m$ as $\cA^e$-modules: $\cA\cong s(\cA)$ is a direct summand of the
free module $\cA^e$, i.e.\ projective. Consequently $Q_0=\cA$, $Q_k=0$ for $k\ge1$, with augmentation the
identity, is a projective resolution of $\cA$. By the comparison theorem (\cref{hochschild:sec:Tor}),
$\HH_k(\cA)$ is the homology of $\cA\otimes_{\cA^e}Q_\bullet$, which vanishes in degrees $k\ge1$. In degree
zero, $\HH_0(\cA)=\cA/[\cA,\cA]=\cA=\C^N$ since $\cA$ is commutative.

(Directly: $\epsilon$ is the idempotent $\sum_je_j\otimes e_j$, and $\ker m$ is generated by $1-\epsilon$.)

%----------------------------------------------------------------------
\subsection*{\Cref{hochschild:ex:matrices}}
Idea: the same argument with the splitting $a\mapsto\sum_jaE_{j1}\otimes E_{1j}$; in degree zero the
commutators are exactly the trace-free matrices.

Let $E_{ij}$ be the matrix units, $E_{ij}E_{kl}=\delta_{jk}E_{il}$. Put $\epsilon=\sum_jE_{j1}\otimes E_{1j}\in\cA\otimes\cA$.
For every $c\in\cA$,
\begin{equation*}
  \sum_jcE_{j1}\otimes E_{1j}=\sum_{i,j}c_{ij}E_{i1}\otimes E_{1j}=\sum_jE_{j1}\otimes E_{1j}c,
\end{equation*}
since $E_{1j}c=\sum_lc_{jl}E_{1l}$ and renaming indices gives the same double sum. Define
$s:\cA\to\cA\otimes\cA=P_0$ by $s(a)=\sum_jaE_{j1}\otimes E_{1j}$. Then
$m(s(a))=\sum_jaE_{j1}E_{1j}=a\sum_jE_{jj}=a$, and, using the identity with $c=b'$,
\begin{equation*}
  (b\otimes b')\cdot s(a)=\sum_jbaE_{j1}\otimes E_{1j}b'=\sum_jbab'E_{j1}\otimes E_{1j}=s(bab'),
\end{equation*}
so $s$ is an $\cA^e$-linear splitting of $m$. As in the previous solution, $\cA$ is a projective
$\cA^e$-module, and $\HH_k(\cA)=0$ for $k\ge1$.

In degree zero $\HH_0(\cA)=\cA/[\cA,\cA]$. The trace vanishes on commutators. Conversely, for $i\ne j$,
$E_{ij}=[E_{ii},E_{ij}]$ and $E_{ii}-E_{jj}=[E_{ij},E_{ji}]$; these $N^2-N+(N-1)=N^2-1$ matrices span
$\ker\tr$ (a matrix with trace zero is a combination of the off-diagonal $E_{ij}$ and of the
$E_{ii}-E_{NN}$). Hence $[\cA,\cA]=\ker\tr$ and $\tr$ induces an isomorphism $\HH_0(M_N(\C))\cong\C$.

%----------------------------------------------------------------------
\subsection*{\Cref{hochschild:ex:delta-low}}
Idea: read off the low degree cases of \cref{hochschild:lem:delta}; $\delta^2=0$ is then the Jacobi
identity.

\emph{$1$-forms.} For $k=1$ the second sum in \cref{hochschild:lem:delta} is empty and
\begin{equation*}
  \delta(a_0\,da_1)=\{a_0,a_1\}.
\end{equation*}
In canonical coordinates a $1$-form is $\alpha=\sum_j(f_j\,dx_j+g_j\,d\xi_j)$, and since
$\{f,x_j\}=\pa_{\xi_j}f$, $\{g,\xi_j\}=-\pa_{x_j}g$,
\begin{equation*}
  \delta\alpha=\sum_{j=1}^n\big(\pa_{\xi_j}f_j-\pa_{x_j}g_j\big).
\end{equation*}
(Consistently, $\iota_\Pi\alpha=0$ and $\delta\alpha=\iota_\Pi d\alpha$; on functions $\delta=0$.)

\emph{$2$-forms.} For $k=2$,
\begin{equation*}
  \delta(a_0\,da_1\wedge da_2)=\{a_0,a_1\}\,da_2-\{a_0,a_2\}\,da_1-a_0\,d\{a_1,a_2\}.
\end{equation*}
For $a_1,a_2$ canonical coordinates $z_a,z_b$ the bracket $\{z_a,z_b\}$ is constant and this becomes
$\delta(f\,dz_a\wedge dz_b)=\{f,z_a\}dz_b-\{f,z_b\}dz_a=\iota_{H_f}(dz_a\wedge dz_b)$, the second statement of
\cref{hochschild:lem:delta}. Every $2$-form is locally a sum of such terms.

\emph{$\delta^2=0$ on $2$-forms.} Since $\delta$ is a differential operator it suffices to check this locally,
on forms $a_0\,da_1\wedge da_2$. By the two formulas,
\begin{equation*}
  \delta^2(a_0\,da_1\wedge da_2)=\{\{a_0,a_1\},a_2\}-\{\{a_0,a_2\},a_1\}-\{a_0,\{a_1,a_2\}\}.
\end{equation*}
The Jacobi identity for the Poisson bracket of $Y$ (the canonical bracket of $T^*X$; see
\cref{wave:ex:poisson}, or check it directly in canonical coordinates) can be written as
$\{a_0,\{a_1,a_2\}\}=\{\{a_0,a_1\},a_2\}+\{a_1,\{a_0,a_2\}\}$ ($\{a_0,\cdot\}=H_{a_0}$ is a derivation of the
bracket), and $\{a_1,\{a_0,a_2\}\}=-\{\{a_0,a_2\},a_1\}$. Hence the right side vanishes, i.e.\ $\delta^2=0$ on
$2$-forms.

%----------------------------------------------------------------------
\subsection*{\Cref{hochschild:ex:circle}}
Idea: $S^*\Sphere^1$ consists of the two directions $\xi>0$, $\xi<0$; each gives a residue trace,
$\int a_{-1}(x,\pm1)\,dx$.

\emph{Cohomology.} With the global coordinate $x\in\R/2\pi\Z$, $T^*\Sphere^1\setminus0=\Sphere^1\times(\R\setminus0)$, and
the dilations act on the second factor, so $S^*\Sphere^1=Y/\R^+=\Sphere^1\times\{+1,-1\}$, the disjoint union of two
circles. Hence $H^0(S^*X)\cong H^1(S^*X)\cong\C^2$ and $H^j(S^*X)=0$ for $j\ge2$. With $n=1$,
\cref{hochschild:thm:main} gives $\HH_k(\cA)\cong H^{2-k}(S^*X)\oplus H^{1-k}(S^*X)$:
\begin{equation*}
  \HH_0\cong0\oplus\C^2=\C^2,\quad\HH_1\cong\C^2\oplus\C^2=\C^4,\quad\HH_2\cong\C^2\oplus0=\C^2,\quad
  \HH_k=0\ (k\ge3).
\end{equation*}

\emph{The two traces.} We use the global coordinate and the left quantization, for which the star product
\eqref{hochschild:eq:star} is $a\star b=\sum_{l\ge0}\frac1{l!}\pa_\xi^la\,D_x^lb$ (all symbols being $2\pi$-periodic in
$x$). For $a=\sum_{m\le M}a_m\in\cA_0$ define
\begin{equation*}
  \tau_\pm(a)=\int_0^{2\pi}a_{-1}(x,\pm1)\,dx .
\end{equation*}
We claim $\tau_\pm(a\star b)=\tau_\pm(b\star a)$. The component of degree $-1$ of $a\star b-b\star a$ is the finite sum
over $l\ge1$ and $p+q-l=-1$ of $\frac1{l!}\big(\pa_\xi^la_p\,D_x^lb_q-\pa_\xi^lb_q\,D_x^la_p\big)$ (the term $l=0$
cancels). Fix such $l,p,q$ and put $c=D_x^la_p$. Integrating by parts on the circle (for fixed $\xi$),
\begin{equation*}
  \int_0^{2\pi}\big(\pa_\xi^la_p\,D_x^lb_q-\pa_\xi^lb_q\,D_x^la_p\big)dx
  =\int_0^{2\pi}\big((-1)^l\pa_\xi^lc\;b_q-c\;\pa_\xi^lb_q\big)dx
  =-\frac{d}{d\xi}\int_0^{2\pi}\sum_{j=0}^{l-1}(-1)^j\pa_\xi^jc\;\pa_\xi^{l-1-j}b_q\,dx,
\end{equation*}
by the identity $u\,\pa^lv-(-1)^l\pa^lu\,v=\pa\big(\sum_{j=0}^{l-1}(-1)^j\pa^ju\,\pa^{l-1-j}v\big)$ (the sum
telescopes). The function of $\xi$ being differentiated is homogeneous of degree $(p-j)+(q-l+1+j)=0$, hence
constant on each half-line $\xi>0$, $\xi<0$, and its derivative vanishes. So $\tau_\pm$ vanishes on
commutators. It also vanishes on $b(\cB^{(1)})$: on each $C_{m,1}$ the map $u\mapsto\tau_\pm(bu)$ is continuous
(only the finitely many $b_l$ with $m-l=-1$ contribute) and vanishes on the dense subspace of finite sums of
products. Hence $\tau_\pm$ are traces on $\cA$, i.e.\ linear forms on $\HH_0(\cA)$.

\emph{Representatives of order $-1$.} Let $A_\pm\in\cA$ be the classes of the operators with symbol
\begin{equation*}
  a_\pm(x,\xi)=\frac1{2\pi|\xi|}\ \text{ for }\pm\xi>0,\qquad a_\pm(x,\xi)=0\ \text{ for }\mp\xi>0,
\end{equation*}
a single homogeneous term of degree $-1$ (smooth on $Y$). On the Fourier modes, $A_+e^{ikx}=\frac{1}{2\pi k}e^{ikx}$ for
$k>0$ and $A_+e^{ikx}=0$ for $k\le0$, modulo smoothing operators; similarly for $A_-$. Then
$\tau_\pm(A_\pm)=1$ and $\tau_\mp(A_\pm)=0$, so $(\tau_+,\tau_-):\HH_0(\cA)\to\C^2$ is surjective, and since
$\dim\HH_0(\cA)=2$ it is an isomorphism. Therefore the classes $[A_+],[A_-]$ form a basis of $\HH_0(\cA)$ consisting
of representatives of order $-1$ (as \cref{hochschild:thm:main} predicts: order $k-n=-1$), every
$A\in\cA$ is congruent modulo $b(\cB^{(1)})$ to $\tau_+(A)A_++\tau_-(A)A_-$, and every trace on $\cA$ is a linear
combination of the two \emph{residue traces} $\tau_+$, $\tau_-$, one for each direction of $\xi$
(\cref{hochschild:rem:residue}).

This agrees with the description of the isomorphism in \cref{hochschild:rem:residue}: for $a$ of degree $-1$,
$a\,\omega=a\,d\xi\wedge dx$, and with $r=|\xi|$ one finds $a\,\omega=\mp a(x,\pm1)\,dx\wedge\frac{dr}r$ on the component
$\pm\xi>0$, so $\alpha''=\mp a(x,\pm1)\,dx$, whose class in $H^1$ of the corresponding circle is
$\mp\tau_\pm(a)$ times the generator dual to $\int_{\Sphere^1}$.

%----------------------------------------------------------------------
\subsection*{\Cref{hochschild:ex:top}}
Idea: in top degree $k=2n$ there are no forms of degree $2n+1$, so $E^2_{n,2n}$ is the kernel of
$\delta$ on homogeneous $2n$-forms, and $\delta(f\,\omega^n)=\iota_{H_f}\omega^n$ vanishes iff $df=0$.

Let $k=2n$. By \cref{hochschild:thm:E2}, $E^2_{p,2n}=0$ unless $p=2n-n=n$, so the hypothesis of
\cref{hochschild:prop:convergence} holds (as it does for all $k$), and every class in $\HH_{2n}(\cA)$ has a
representative cycle $c\in F_n\cB^{(2n)}$, i.e.\ of order at most $n$; moreover $\Theta[c]=[c_n]\in E^2_{n,2n}$ is an
isomorphism, $c_n\in C_{n,2n}$ being a $b_0$-cycle whose class in $E^1_{n,2n}\cong\Omega^{2n}_{(n)}(Y)$ is
$\alpha=\varepsilon(c_n)$ (\eqref{hochschild:eq:E1}).

We identify $E^2_{n,2n}$. By \cref{hochschild:prop:d1} it is the homology at $\Omega^{2n}_{(n)}(Y)$ of
$\Omega^{2n+1}_{(n+1)}\xrightarrow{\delta}\Omega^{2n}_{(n)}\xrightarrow{\delta}\Omega^{2n-1}_{(n-1)}$. Since $\dim Y=2n$,
$\Omega^{2n+1}=0$, so $E^2_{n,2n}=\{\alpha\in\Omega^{2n}_{(n)}(Y);\ \delta\alpha=0\}$. Every $2n$-form is
$\alpha=f\,\mathrm{vol}$ with $\mathrm{vol}=\omega^n/n!$ and $f\in\CI(Y)$; as $\mathrm{vol}$ is homogeneous of degree $n$
($\rho_t^*\omega=t\omega$), $\alpha\in\Omega^{2n}_{(n)}$ iff $f$ is homogeneous of degree $0$. Since $\mathrm{vol}$ has constant
coefficients in canonical coordinates, the second statement of \cref{hochschild:lem:delta} gives
\begin{equation*}
  \delta(f\,\mathrm{vol})=\iota_{H_f}\mathrm{vol}.
\end{equation*}
As $\mathrm{vol}$ is a volume form, $\iota_v\mathrm{vol}=0$ only for $v=0$; and $H_f=0$ iff $df=0$ (the map $df\mapsto H_f$
is invertible, $\omega$ being non-degenerate). Hence $\delta\alpha=0$ iff $f$ is locally constant on $Y$ (such $f$ are
automatically homogeneous of degree $0$). Therefore
\begin{equation*}
  E^2_{n,2n}=\{f\,\omega^n;\ f\text{ locally constant on }Y\}\cong H^0(Y)=H^0(S^*X),
\end{equation*}
in agreement with \cref{hochschild:thm:E2} ($H^{2n-k}\oplus H^{2n-k-1}=H^0(S^*X)\oplus0$ for $k=2n$). The same
conclusion follows from \cref{hochschild:prop:star-d}: $*$ maps $\Omega^{2n}_{(n)}$ onto $\Omega^0_{(0)}$, $*\mathrm{vol}$ is the
non-zero constant $\Lambda(\mathrm{vol},\mathrm{vol})$, and $\delta=-*^{-1}d*$ on $2n$-forms, so $\delta(f\,\mathrm{vol})=0$ iff
$df=0$.

Putting things together: every class in $\HH_{2n}(\cA)$ is represented by a cycle $c$ of order at most $n$ whose leading
part $c_n$ satisfies $\varepsilon(c_n)=f\,\omega^n$ with $f$ locally constant on $Y$ (we may absorb $1/n!$ into $f$), and
$[c]\mapsto f$ is an isomorphism of $\HH_{2n}(\cA)$ onto the locally constant functions on $Y$. Conversely,
given such an $f$, \cref{hochschild:thm:commutative} provides a $b_0$-cycle $x\in C_{n,2n}$ with
$\varepsilon(x)=f\,\omega^n$, and the surjectivity part of the proof of \cref{hochschild:prop:convergence} corrects
$x$ by chains of order $<n$ to a cycle $c$ with $c_n=x$.
